\documentclass[11pt,letterpaper]{article}

%% ── Packages ────────────────────────────────────────────────────────────────
\usepackage[margin=1in]{geometry}
\usepackage{amsmath,amssymb,amsthm}
\usepackage{mathtools}
\usepackage{booktabs}
\usepackage{longtable}
\usepackage{array}
\usepackage{xcolor}
\usepackage{hyperref}
\usepackage{listings}
\usepackage{fancyhdr}
\usepackage{titlesec}
\usepackage{enumitem}
\usepackage{mdframed}
\usepackage{graphicx}
\usepackage{microtype}
\usepackage{cite}
\usepackage{datetime2}

%% ── Colours ─────────────────────────────────────────────────────────────────
\definecolor{mfblue}{RGB}{0,74,173}
\definecolor{mfgray}{RGB}{90,90,90}
\definecolor{codebg}{RGB}{248,248,248}
\definecolor{codeframe}{RGB}{180,180,180}
\definecolor{rustred}{RGB}{165,42,42}
\definecolor{tsblue}{RGB}{0,112,186}

%% ── Hyperref ─────────────────────────────────────────────────────────────────
\hypersetup{
  colorlinks=true,
  linkcolor=mfblue,
  urlcolor=mfblue,
  citecolor=mfgray,
  pdftitle={Multiplicity-Crypto: WASM Pedersen Commitment Integration and QSBS Architecture},
  pdfauthor={Ryan Van Gelder, Multiplicity Foundation},
}

%% ── Headers ─────────────────────────────────────────────────────────────────
\pagestyle{fancy}
\fancyhf{}
\fancyhead[L]{\small\textcolor{mfgray}{Multiplicity Foundation --- Defensive Publication}}
\fancyhead[R]{\small\textcolor{mfgray}{April 2026}}
\fancyfoot[C]{\small\thepage}
\renewcommand{\headrulewidth}{0.4pt}

%% ── Section styling ──────────────────────────────────────────────────────────
\titleformat{\section}{\large\bfseries\color{mfblue}}{}{0em}{}[\titlerule]
\titleformat{\subsection}{\normalsize\bfseries\color{mfblue}}{}{0em}{}

%% ── Theorem environments ─────────────────────────────────────────────────────
\newtheorem{definition}{Definition}[section]
\newtheorem{theorem}{Theorem}[section]
\newtheorem{proposition}{Proposition}[section]
\newtheorem{remark}{Remark}[section]

%% ── Code listings ────────────────────────────────────────────────────────────
\lstset{
  backgroundcolor=\color{codebg},
  frame=single,
  rulecolor=\color{codeframe},
  basicstyle=\ttfamily\small,
  keywordstyle=\color{mfblue}\bfseries,
  stringstyle=\color{rustred},
  commentstyle=\color{mfgray}\itshape,
  breaklines=true,
  numbers=left,
  numberstyle=\tiny\color{mfgray},
  numbersep=6pt,
  tabsize=2,
  captionpos=b,
}

\lstdefinelanguage{Rust}{
  keywords={use,pub,fn,let,mut,struct,impl,trait,return,if,else,match,for,
             while,loop,type,const,static,extern,mod,crate,self,super,where,
             as,in,ref,true,false,String,bool,u8,u32,u64,usize,i32,i64,Vec},
  morecomment=[l]{//},
  morecomment=[s]{/*}{*/},
  morestring=[b]",
  sensitive=true,
}

\lstdefinelanguage{TypeScript}{
  keywords={import,export,from,const,let,var,function,async,await,return,
             if,else,try,catch,throw,new,class,interface,type,extends,
             implements,true,false,null,undefined,string,number,boolean,
             Promise,void},
  morecomment=[l]{//},
  morecomment=[s]{/*}{*/},
  morestring=[b]",
  morestring=[b]',
  sensitive=true,
}

%% ─────────────────────────────────────────────────────────────────────────────
\begin{document}

%% ── Title page ───────────────────────────────────────────────────────────────
\begin{titlepage}
\centering
\vspace*{2cm}
{\Huge\bfseries\color{mfblue}
  Multiplicity-Crypto:\\[0.4em]
  WASM Pedersen Commitment Integration\\[0.3em]
  and Quantum-State-Based Security Architecture
\par}
\vspace{1.2cm}
{\large\textbf{Defensive Publication \& Prior Art Disclosure}\par}
\vspace{0.4cm}
{\large\textcolor{mfgray}{Technical Report v1.0.0}\par}
\vspace{1.4cm}
\hrule height 0.8pt
\vspace{0.6cm}
{\large
  \textbf{Ryan Van Gelder}\\[0.2em]
  \textcolor{mfgray}{Founder \& Principal Researcher}\\[0.2em]
  Multiplicity Foundation (501(c)(3))\\[0.2em]
  \texttt{MultiplicityFoundation/PhaseMirror-HQ}
\par}
\vspace{0.6cm}
\hrule height 0.4pt
\vspace{1.6cm}
{\normalsize
\begin{tabular}{ll}
  \textbf{Document Date:} & April 25, 2026 \\
  \textbf{Repository:}    & \texttt{packages/multiplicity-crypto} \\
  \textbf{License:}       & MIT \& CC BY-SA 4.0 \\
  \textbf{Prior Art:}     & Public Disclosure for Defensive Patent Purpose \\
\end{tabular}
\par}
\vspace{2cm}
\begin{mdframed}[linecolor=mfblue,linewidth=1.2pt,innertopmargin=8pt,innerbottommargin=8pt]
\small\textit{This document constitutes a defensive publication establishing prior art
for the cryptographic architectures, mathematical frameworks, and software
implementations described herein. Publication date: April 25, 2026. All rights
reserved under the Multiplicity Prime Materia License except where otherwise noted.
This publication is intended to place these inventions irrevocably in the public
domain for defensive purposes, preventing third-party patent claims on the described
techniques.}
\end{mdframed}
\end{titlepage}

%% ── Abstract ─────────────────────────────────────────────────────────────────
\begin{abstract}
\noindent
We present the design, mathematical foundations, and implementation roadmap for
\textbf{multiplicity-crypto}, a cryptographic primitive library operating as a core
component of the Phase Mirror governance protocol system. This report documents the
architectural decisions, prior art, and technical implementation plan for three
integrated subsystems: (1) a native WebAssembly (WASM) Pedersen commitment engine
over the BN254 elliptic curve, suitable for direct ZK-SNARK circuit verification under
Groth16 and Plonk backends; (2) a Quantum-State-Based Security (QSBS) encryption
engine unifying classical and quantum data through a frequency-mapped, multiplicity-
indexed encryption operator; and (3) a QKD Hybrid Encryption transport protocol
(v1.0.1) providing transcript-bound, HKDF-derived directional AEAD keys with
optional multiplicity profile extensions. The document establishes architectural
layer boundaries, mathematical formalism, code structure, known-vector test
methodology, and a sequenced implementation roadmap. All content constitutes
original work by Ryan Van Gelder and the Multiplicity Foundation.
\end{abstract}

\tableofcontents
\newpage

%% ═════════════════════════════════════════════════════════════════════════════
\section{Executive Summary}
%% ═════════════════════════════════════════════════════════════════════════════

The \texttt{multiplicity-crypto} package (\texttt{@phase-mirror/multiplicity-crypto})
is a TypeScript/Rust/WASM cryptographic primitive library serving as the lowest
cryptographic layer of the Phase Mirror governance protocol stack. As of April 2026,
the package contains five source modules:
\texttt{commitment.ts}, \texttt{merkle.ts}, \texttt{bls.ts}, \texttt{randomness.ts},
and \texttt{index.ts}, backed by a Rust crate
(\texttt{multiplicity-crypto-wasm}) that implements genuine elliptic-curve Pedersen
commitments via \texttt{wasm-bindgen}.

\subsection{Current Completion State}

\begin{longtable}{@{}lllp{5.2cm}@{}}
\toprule
\textbf{Component} & \textbf{Status} & \textbf{Implementation} & \textbf{Blocker / Notes} \\
\midrule
\endhead
Commitment (Pedersen) & Placeholder & SHA-256 fallback & WASM build not yet executed; Rust code complete but on wrong curve (secp256k1) \\
Merkle Root & Placeholder & SHA-256 & Depends on Pedersen WASM landing \\
BLS12-381 & Stub & SHA-256 stub & Reserved for signature aggregation \\
Randomness & Production & HMAC-SHA256 & No test coverage despite production label \\
Poseidon Hash & Deprecated & SHA-256 stub & Dead code; files not yet deleted \\
WASM Bindings & Missing & N/A & \texttt{src/wasm/} directory absent \\
CI Pipeline & Missing & N/A & No automated build or test \\
QSBS Engine & Unimplemented & N/A & Mathematical spec complete; no code \\
QKD Adapter & Unimplemented & N/A & Simulator backend planned \\
\bottomrule
\caption{Component completion matrix as of April 25, 2026}
\end{longtable}

\subsection{Key Architectural Decisions Established}

\begin{enumerate}[leftmargin=*]
  \item \textbf{Curve selection:} BN254 (\texttt{alt\_bn128}) is mandated for all
        Pedersen commitments to ensure native Groth16/Plonk circuit verifiability.
        The existing \texttt{k256} (secp256k1) dependency must be replaced with
        \texttt{ark-bn254}.
  \item \textbf{WASM build pipeline:} \texttt{wasm-pack build --target bundler}
        must generate \texttt{src/wasm/multiplicity\_crypto\_wasm.\{js,d.ts,wasm\}}
        and be automated in CI.
  \item \textbf{QSBS integration:} The Quantum-State-Based Security engine operates
        at Layer 3, consuming BN254 commitment outputs as its hash handle \(h\).
  \item \textbf{QKD transport:} The v1.0.1 hybrid QKD protocol provides the key
        substrate \(K_t\), with multiplicity operator state \(M_t\) embedded in
        HKDF \texttt{info} strings.
  \item \textbf{Implementation order:} Architecture doc \(\to\) BN254 Rust/WASM \(\to\)
        transcript chain \(\to\) HKDF key derivation \(\to\) frequency mapping \(\to\)
        AEAD \(\to\) QKD simulator \(\to\) multiplicity/feedback dynamics \(\to\)
        KAT runner.
\end{enumerate}

%% ═════════════════════════════════════════════════════════════════════════════
\section{Mathematical Foundations}
%% ═════════════════════════════════════════════════════════════════════════════

\subsection{Pedersen Commitments over BN254}

\begin{definition}[BN254 Pedersen Commitment]
Let \(\mathbb{G}_1\) be the prime-order subgroup of the BN254 elliptic curve with
generator \(G\). Let \(\mathbb{F}_r\) denote the scalar field of order \(r\).
Define an independent generator \(H \in \mathbb{G}_1\) via:
\[
  H = \left( \text{SHA-256}(\text{``pedersen\_h''} \;\|\; G_{\text{enc}}) \bmod r \right) \cdot G
\]
For value \(v \in \mathbb{F}_r\) and blinding factor \(\rho \in \mathbb{F}_r\), the
Pedersen commitment is:
\[
  \mathbf{C}(v, \rho) = \rho \cdot G + v \cdot H \;\in\; \mathbb{G}_1 \tag{1}
\]
The commitment is binding under the discrete logarithm assumption and perfectly
hiding.
\end{definition}

\begin{definition}[Scalar Derivation]
Scalars are derived from string inputs via domain-separated SHA-256:
\[
  \text{derive\_scalar}(d, s) = \text{SHA-256}(d \;\|\; s) \bmod r \;\in\; \mathbb{F}_r \tag{2}
\]
where \(d\) is a domain string (\texttt{"pedersen\_value"} or
\texttt{"pedersen\_blinding"}) and \(s\) is the input string.
\end{definition}

\begin{proposition}[Circuit Verifiability]
Commitments of the form \(\mathbf{C}(v,\rho)\) over BN254 are natively verifiable
inside Groth16 and Plonk arithmetic circuits without custom gate overhead, because
BN254 is the native curve of the \texttt{bn256} pairing used in both proof systems.
By contrast, secp256k1 (k256) commitments require \(O(256)\) non-native field
arithmetic constraints per verification, rendering them impractical for circuit use.
\end{proposition}

\subsection{Multiplicity Theory Integration}

\begin{definition}[Multiplicity Operator]
Let \(\mathcal{P} = \{p_1, p_2, \ldots\}\) denote the ordered sequence of primes.
The multiplicity operator \(M_t\) at time \(t\) is a function:
\[
  M_t : \mathcal{S} \to \mathbb{Z}_{>0}^k, \quad
  M_t(\sigma) = \bigoplus_{i=1}^{k} \nu_{p_i}(\sigma) \tag{3}
\]
where \(\nu_{p_i}(\sigma)\) denotes the \(p_i\)-adic valuation of the system state
encoding \(\sigma\), and \(\oplus\) denotes prime-indexed tensor composition.
\end{definition}

\begin{definition}[Coupling Tensor]
The coupling tensor \(T_t \in \mathbb{R}^{n \times m}\) describes interactions
between \(n\) classical subsystem components and \(m\) quantum amplitude modes:
\[
  (T_t)_{ij} = \langle c_i \,|\, q_j \rangle_t \tag{4}
\]
where \(c_i\) are classical frequency components and \(q_j\) are quantum amplitude
frequency components at time \(t\).
\end{definition}

\subsection{Unified System State}

\begin{definition}[QSBS System State]
The full system state at discrete time \(t\) is the quintuple:
\[
  S_t = \bigl(x_t,\; |\psi_t\rangle,\; \sigma_t,\; M_t,\; T_t\bigr) \tag{5}
\]
where:
\begin{itemize}
  \item \(x_t \in \{0,1\}^n\) is the classical payload bitstring
  \item \(|\psi_t\rangle = \sum_{j=1}^{m} \alpha_j |j\rangle\) is the quantum state
        with complex amplitudes \(\alpha_j\)
  \item \(\sigma_t\) is the protocol transcript (sequence numbers, labels, message history)
  \item \(M_t\) is the multiplicity operator state (Definition 2.3)
  \item \(T_t\) is the coupling tensor (Definition 2.4)
\end{itemize}
\end{definition}

\subsection{Frequency Mapping}

\begin{definition}[Classical Frequency Map]
The classical frequency map \(F_c : \{0,1\}^n \to \mathbb{R}^n\) is defined by:
\[
  F_c(x_i) = \begin{cases} f_{c,i} & \text{if } x_i = 0 \\ f'_{c,i} & \text{if } x_i = 1 \end{cases} \tag{6}
\]
\end{definition}

\begin{definition}[Quantum Frequency Map]
The quantum frequency map \(F_q : \mathbb{C}^m \to \mathbb{R}^m\) encodes both
magnitude and phase of each amplitude:
\[
  F_q(\alpha_j) = f_{q,j}, \qquad f_{q,j} = \bigl(|\alpha_j|, \arg(\alpha_j)\bigr) \tag{7}
\]
\end{definition}

\begin{definition}[Combined Frequency Set]
\[
  \mathbf{F}_t = F_c(x_t) \cup F_q\bigl(|\psi_t\rangle\bigr) \tag{8}
\]
\end{definition}

\subsection{Encryption Operator}

\begin{definition}[QSBS Encryption Operator]
The unified encryption operator is:
\[
  E(t) = M_t \cdot (T_t \cdot S_t) + F(S_t) \tag{9}
\]
where \(F : \mathcal{S} \to \mathcal{S}\) is the non-linear feedback function
(Definition 2.6 below). In the wire-protocol instantiation:
\[
  C_t = \text{AEAD}_{K_{\text{enc},d,t}}\!\bigl(\text{nonce}_{d,t},\; x_t;\;
        \text{ad} = \sigma_t \,\|\, \mathbf{F}_t \,\|\, M_t \,\|\, T_t \bigr) \tag{10}
\]
\end{definition}

\subsection{Key Derivation with Multiplicity Profile}

\begin{definition}[Transcript Context Hash]
The per-sender transcript hash chain follows SHA-256 with concatenated per-sender
chains combined as:
\[
  \text{contexthash}_t = H_{\text{ctx}}(\sigma_t) = \text{SHA-256}
    \bigl(\text{chain}_A \,\|\, \text{chain}_B\bigr) \tag{11}
\]
\end{definition}

\begin{definition}[HKDF Key Derivation with Multiplicity]
Directional AEAD keys are derived as:
\[
  K_{\text{enc},d,t} = \text{HKDF}\!\left(K_t;\; \text{salt}_t,\;
    \text{info}_{d,\,\text{contexthash}_t,\,M_t}\right) \tag{12}
\]
where \(d \in \{\text{A2B}, \text{B2A}\}\) and the \texttt{info} string is:
\[
  \text{info}_{d,h,M} = \texttt{"QSBS-v1:enc:"} \,\|\, d \,\|\, h \,\|\, M_{\text{label}} \tag{13}
\]
\(\,M_{\text{label}}\) is a fixed-length (8-byte) encoding of the current multiplicity
state index, embedded without changing ciphertext semantics.
\end{definition}

\subsection{Feedback Dynamics}

\begin{definition}[Non-Linear Feedback Function]
The multiplicity and coupling states evolve according to:
\[
  \bigl(M_{t+1}, T_{t+1}\bigr) = f\!\bigl(M_t, T_t;\;
    \varepsilon_t, \lambda_t, \text{QBER}_t, \ell_t\bigr) \tag{14}
\]
where \(\varepsilon_t\) is the error rate, \(\lambda_t\) is observed latency,
\(\text{QBER}_t\) is the quantum bit error rate, and \(\ell_t\) is the system load.
\end{definition}

\begin{remark}
Under benign conditions, \(f\) converges to stable multiplicity profiles that
maximize throughput while maintaining transcript binding. Under elevated error or
attack conditions, \(f\) increases key rotation rate and coupling tensor weight,
observable as increased latency, key-pool consumption, and multiplicity index
variance.
\end{remark}

%% ═════════════════════════════════════════════════════════════════════════════
\section{Architecture and Layer Model}
%% ═════════════════════════════════════════════════════════════════════════════

\subsection{Layer Stack}

\begin{longtable}{@{}cllll@{}}
\toprule
\textbf{Layer} & \textbf{Symbol} & \textbf{Description} & \textbf{Target File} & \textbf{Status} \\
\midrule
\endhead
L0 & --- & BN254 curve primitives & \texttt{rust/src/lib.rs} & Needs curve swap \\
L1 & \(h\) & WASM commitment API & \texttt{src/commitment.ts} & Blocked on L0 \\
L2 & \(\mathbf{F}_t\) & Frequency mapping & \texttt{src/frequency.ts} & Not created \\
L3 & \(\text{ctx}_t\) & Transcript hash chain & \texttt{src/transcript.ts} & Not created \\
L4 & \(K_{\text{enc}}\) & HKDF key derivation & \texttt{src/keyderivation.ts} & Not created \\
L5 & \(M_t\) & Multiplicity operator & \texttt{src/multiplicity.ts} & Not created \\
L6 & \(f(\cdot)\) & Feedback dynamics & \texttt{src/feedback.ts} & Not created \\
L7 & \(C_t\) & AEAD encrypt/decrypt & \texttt{src/aead.ts} & Not created \\
L8 & \(K_t\) & QKD adapter/simulator & \texttt{src/qkd.ts} & Not created \\
L9 & --- & KAT runner (Tests 1--16) & \texttt{src/kat.test.ts} & Not created \\
\bottomrule
\caption{Full architecture layer stack with file targets}
\end{longtable}

\subsection{Dependency Graph}

\[
  \texttt{lib.rs} \;\to\; \texttt{commitment.ts} \;\to\; \texttt{transcript.ts}
  \;\to\; \texttt{keyderivation.ts} \;\to\; \texttt{aead.ts}
  \;\to\; \texttt{qkd.ts}
\]
\[
  \texttt{frequency.ts} \;\to\; \texttt{aead.ts}
\]
\[
  \texttt{multiplicity.ts} + \texttt{feedback.ts} \;\to\; \texttt{keyderivation.ts}
\]

%% ═════════════════════════════════════════════════════════════════════════════
\section{Code: Rust BN254 Pedersen Implementation}
%% ═════════════════════════════════════════════════════════════════════════════

\subsection{Cargo.toml (Required)}

The existing \texttt{k256} dependency must be replaced with the \texttt{arkworks}
BN254 stack. The \texttt{rust/Cargo.toml} target state is:

\begin{lstlisting}[language={},caption={packages/multiplicity-crypto/rust/Cargo.toml (target state)}]
[package]
name    = "multiplicity-crypto-wasm"
version = "0.2.0"
edition = "2021"

[lib]
crate-type = ["cdylib"]

[dependencies]
wasm-bindgen = "0.2"
ark-bn254     = "0.4"
ark-ec        = { version = "0.4", features = ["curve"] }
ark-ff        = "0.4"
ark-std       = "0.4"
sha2          = "0.10"
hex           = "0.4"
getrandom     = { version = "0.2", features = ["js"] }

[dev-dependencies]
wasm-bindgen-test = "0.3"
\end{lstlisting}

\subsection{lib.rs (BN254 Pedersen)}

\begin{lstlisting}[language=Rust,caption={packages/multiplicity-crypto/rust/src/lib.rs (target state)}]
use ark_bn254::{G1Projective, Fr};
use ark_ec::Group;
use ark_ff::{PrimeField, BigInteger};
use sha2::{Sha256, Digest};
use wasm_bindgen::prelude::*;

/// Derive a BN254 scalar from a domain-separated SHA-256 hash.
/// Implements equation (2): derive_scalar(d, s) = SHA-256(d || s) mod r
fn derive_scalar(domain: &str, data: &str) -> Fr {
    let mut h = Sha256::new();
    h.update(domain.as_bytes());
    h.update(data.as_bytes());
    let bytes = h.finalize();
    Fr::from_be_bytes_mod_order(&bytes)
}

/// Derive the independent Pedersen generator H.
/// H = SHA-256("pedersen_h" || G_enc) * G  -- see equation (1)
fn pedersen_generator_h() -> G1Projective {
    let g = G1Projective::generator();
    let h_scalar = derive_scalar("pedersen_h", "generator");
    g * h_scalar
}

/// Compute C(v, rho) = rho*G + v*H over BN254.
/// Implements equation (1): C(v, rho) = rho*G + v*H
/// Returns the x-coordinate of the affine commitment point as 0x-prefixed hex.
#[wasm_bindgen]
pub fn compute_pedersen_commitment(value: &str, salt: &str) -> String {
    let g = G1Projective::generator();
    let h = pedersen_generator_h();
    let v   = derive_scalar("pedersen_value",    value);
    let rho = derive_scalar("pedersen_blinding", salt);
    let commitment = (g * rho + h * v).into_affine();
    let x_bytes = commitment.x.into_bigint().to_bytes_be();
    format!("0x{}", hex::encode(&x_bytes))
}

/// Verify a Pedersen commitment deterministically.
#[wasm_bindgen]
pub fn verify_pedersen_commitment(
    commitment: &str,
    value: &str,
    salt: &str,
) -> bool {
    compute_pedersen_commitment(value, salt) == commitment
}

#[cfg(test)]
mod tests {
    use super::*;

    // Known-vector tests -- values must be pre-computed and pinned.
    #[test]
    fn test_commitment_deterministic() {
        let c1 = compute_pedersen_commitment("leaf", "salt");
        let c2 = compute_pedersen_commitment("leaf", "salt");
        assert_eq!(c1, c2);
    }

    #[test]
    fn test_commitment_0x_prefix_64_bytes() {
        let c = compute_pedersen_commitment("test", "nonce");
        assert!(c.starts_with("0x"));
        assert_eq!(c.len(), 66); // "0x" + 64 hex chars = 32 bytes
    }

    #[test]
    fn test_commitment_binding() {
        let c = compute_pedersen_commitment("secret", "salt");
        assert!(verify_pedersen_commitment(&c, "secret", "salt"));
        assert!(!verify_pedersen_commitment(&c, "wrong", "salt"));
    }
}
\end{lstlisting}

%% ═════════════════════════════════════════════════════════════════════════════
\section{Code: TypeScript WASM Integration}
%% ═════════════════════════════════════════════════════════════════════════════

\subsection{commitment.ts}

\begin{lstlisting}[language=TypeScript,caption={packages/multiplicity-crypto/src/commitment.ts}]
import { createHash } from 'crypto';
import init, {
  compute_pedersen_commitment,
  verify_pedersen_commitment,
} from './wasm/multiplicity_crypto_wasm';

let wasmInitialized = false;

async function ensureWasmInitialized(): Promise<void> {
  if (!wasmInitialized) {
    await init();
    wasmInitialized = true;
  }
}

/**
 * Compute a BN254 Pedersen commitment: C(v, rho) = rho*G + v*H
 * Falls back to SHA-256 only if WASM is unavailable (test/CI environments).
 * WARNING: SHA-256 fallback is NOT circuit-verifiable. Do not use in production.
 */
export async function computeCommitment(
  leaf: string,
  salt: string,
): Promise<string> {
  try {
    await ensureWasmInitialized();
    return compute_pedersen_commitment(leaf, salt);
  } catch (error) {
    // FALLBACK: SHA-256 placeholder -- NOT BN254, NOT circuit-verifiable
    console.warn('[multiplicity-crypto] WASM unavailable: SHA-256 fallback active.');
    const digest = createHash('sha256')
      .update(leaf, 'utf8')
      .update(':commit:', 'utf8')
      .update(salt, 'utf8')
      .digest('hex');
    return `0x${digest}`;
  }
}

export async function verifyCommitment(
  commitment: string,
  leaf: string,
  salt: string,
): Promise<boolean> {
  try {
    await ensureWasmInitialized();
    return verify_pedersen_commitment(commitment, leaf, salt);
  } catch {
    const computed = createHash('sha256')
      .update(leaf, 'utf8')
      .update(':commit:', 'utf8')
      .update(salt, 'utf8')
      .digest('hex');
    return commitment === `0x${computed}`;
  }
}
\end{lstlisting}

\subsection{transcript.ts (New Module)}

\begin{lstlisting}[language=TypeScript,caption={packages/multiplicity-crypto/src/transcript.ts}]
import { createHash } from 'crypto';

/**
 * Per-sender SHA-256 transcript hash chain.
 * Implements equation (11): contexthash_t = SHA-256(chain_A || chain_B)
 */
export class TranscriptChain {
  private chain: Buffer;

  constructor() {
    this.chain = Buffer.alloc(32, 0); // initial zero hash
  }

  update(message: Buffer): void {
    this.chain = createHash('sha256')
      .update(this.chain)
      .update(message)
      .digest();
  }

  getChain(): Buffer { return Buffer.from(this.chain); }
}

export function computeContextHash(
  chainA: Buffer,
  chainB: Buffer,
): string {
  const hash = createHash('sha256')
    .update(chainA)
    .update(chainB)
    .digest('hex');
  return `0x${hash}`;
}
\end{lstlisting}

\subsection{keyderivation.ts (New Module)}

\begin{lstlisting}[language=TypeScript,caption={packages/multiplicity-crypto/src/keyderivation.ts}]
import { createHmac } from 'crypto';

type Direction = 'A2B' | 'B2A';

/**
 * HKDF-Extract: PRK = HMAC-SHA256(salt, IKM)
 */
function hkdfExtract(salt: Buffer, ikm: Buffer): Buffer {
  return createHmac('sha256', salt).update(ikm).digest();
}

/**
 * HKDF-Expand: OKM = T(1) where T(i) = HMAC-SHA256(PRK, T(i-1) || info || i)
 */
function hkdfExpand(prk: Buffer, info: Buffer, length: number): Buffer {
  const okm = Buffer.alloc(length);
  let t = Buffer.alloc(0);
  let offset = 0;
  for (let i = 1; offset < length; i++) {
    t = createHmac('sha256', prk)
      .update(t).update(info).update(Buffer.from([i]))
      .digest();
    t.copy(okm, offset);
    offset += t.length;
  }
  return okm.slice(0, length);
}

/**
 * Derive directional AEAD key with optional multiplicity label M_t.
 * Implements equation (12)-(13):
 * K_enc,d,t = HKDF(K_t; salt_t, "QSBS-v1:enc:" || d || contexthash || M_label)
 */
export function deriveEncryptionKey(
  sharedKey: Buffer,
  salt: Buffer,
  direction: Direction,
  contextHash: string,
  multiplicityLabel?: Buffer,  // 8-byte M_t encoding; equation (13)
): Buffer {
  const mLabel = multiplicityLabel ?? Buffer.alloc(8, 0);
  const info = Buffer.concat([
    Buffer.from('QSBS-v1:enc:'),
    Buffer.from(direction),
    Buffer.from(contextHash),
    mLabel,
  ]);
  const prk = hkdfExtract(salt, sharedKey);
  return hkdfExpand(prk, info, 32);
}
\end{lstlisting}

%% ═════════════════════════════════════════════════════════════════════════════
\section{Code: QSBS Engine Scaffold}
%% ═════════════════════════════════════════════════════════════════════════════

\begin{lstlisting}[language=TypeScript,caption={packages/multiplicity-crypto/src/qsbs.ts (scaffold)}]
import { computeCommitment } from './commitment';

/**
 * Quantum-State-Based Security encryption operator.
 * Implements equation (9): E(t) = M(t) * (T * S) + F(S)
 *
 * At this scaffold stage, M, T, and F are represented as
 * numeric placeholders pending full multiplicity-operator implementation.
 */
export interface QSBSState {
  classicalData: string;        // x_t
  quantumAmplitudes: number[];  // alpha_j for |psi_t>
  multiplicityIndex: number;    // scalar projection of M_t
  couplingWeight: number;       // scalar projection of T_t
}

/**
 * Compute frequency mapping F = F_c(x) union F_q(psi)
 * Implements equations (6), (7), (8).
 */
export function computeFrequencyMapping(state: QSBSState): number[] {
  // F_c: map each classical bit to a frequency bin
  const classicalFreqs = state.classicalData
    .split('')
    .map((c, i) => (c === '0' ? i * 2.0 : i * 2.0 + 1.0));

  // F_q: map each amplitude to (magnitude, phase) pair
  const quantumFreqs = state.quantumAmplitudes.flatMap((a, j) => [
    Math.abs(a),
    Math.atan2(0, a), // phase; extend for complex amplitudes
  ]);

  return [...classicalFreqs, ...quantumFreqs];
}

/**
 * QSBS encryption: produces a commitment handle h = H(F_t)
 * that serves as the encryption handle for the AEAD layer.
 * Implements equation (9) in hash-commitment form.
 */
export async function qsbsEncryptionHandle(
  state: QSBSState,
  salt: string,
): Promise<string> {
  const freqs = computeFrequencyMapping(state);
  const freqString = freqs.map(f => f.toFixed(8)).join(':');
  const saltWithMultiplicity = `${salt}:M=${state.multiplicityIndex}`;
  return computeCommitment(freqString, saltWithMultiplicity);
}
\end{lstlisting}

%% ═════════════════════════════════════════════════════════════════════════════
\section{CI Pipeline}
%% ═════════════════════════════════════════════════════════════════════════════

\begin{lstlisting}[language={},caption={.github/workflows/multiplicity-crypto.yml}]
name: multiplicity-crypto

on:
  push:
    paths: ['packages/multiplicity-crypto/**']
  pull_request:
    paths: ['packages/multiplicity-crypto/**']

jobs:
  build-and-test:
    runs-on: ubuntu-latest
    steps:
      - uses: actions/checkout@v4

      - name: Install Rust toolchain
        uses: dtolnay/rust-toolchain@stable
        with:
          targets: wasm32-unknown-unknown

      - name: Install wasm-pack
        run: curl https://rustwasm.github.io/wasm-pack/installer/init.sh -sSf | sh

      - name: Build WASM (BN254 Pedersen)
        working-directory: packages/multiplicity-crypto/rust
        run: |
          wasm-pack build --target bundler \
            --out-dir ../src/wasm \
            --out-name multiplicity_crypto_wasm

      - name: Install Node dependencies
        run: pnpm install --frozen-lockfile

      - name: Run Rust unit tests
        working-directory: packages/multiplicity-crypto/rust
        run: cargo test

      - name: Run TypeScript tests
        run: pnpm --filter @phase-mirror/multiplicity-crypto test

      - name: Assert no WASM fallback in test output
        run: |
          pnpm --filter @phase-mirror/multiplicity-crypto test 2>&1 | \
          grep -q "WASM unavailable" && exit 1 || exit 0
\end{lstlisting}

%% ═════════════════════════════════════════════════════════════════════════════
\section{Implementation Roadmap}
%% ═════════════════════════════════════════════════════════════════════════════

\begin{longtable}{@{}clp{5.5cm}lc@{}}
\toprule
\textbf{Step} & \textbf{Owner File} & \textbf{Deliverable} & \textbf{Horizon} & \textbf{Gate} \\
\midrule
\endhead
1 & \texttt{ARCHITECTURE.md} & Layer stack, curve decision, input/output contracts & 1 day & --- \\
2 & \texttt{rust/Cargo.toml} & Swap \texttt{k256} $\to$ \texttt{ark-bn254} & 1 day & --- \\
3 & \texttt{rust/src/lib.rs} & BN254 Pedersen; 3 known-vector \texttt{\#[test]} & 2 days & Gate 2 \\
4 & CI & WASM build step; \texttt{dist/} + \texttt{target/} in \texttt{.gitignore} & 2 days & --- \\
5 & \texttt{src/wasm/} & \texttt{wasm-pack build} output committed or CI-generated & 2 days & --- \\
6 & \texttt{src/transcript.ts} & SHA-256 chain, \texttt{computeContextHash}; Tests 1--3 & 3 days & --- \\
7 & \texttt{src/keyderivation.ts} & HKDF + $M_t$ label slot; HKDF test vectors & 3 days & --- \\
8 & \texttt{src/frequency.ts} & $F_c$ and $F_q$ mappings; unit tests & 4 days & --- \\
9 & \texttt{src/aead.ts} & AEAD encrypt/decrypt with \texttt{contexthash} AD & 5 days & Gate 3 \\
10 & \texttt{src/qkd.ts} & HKDF-based QKD simulator; \texttt{SIMULATED\_QKD} label & 1 week & --- \\
11 & \texttt{src/multiplicity.ts} & $M_t$ operator; prime-indexed state encoding & 2 weeks & --- \\
12 & \texttt{src/feedback.ts} & $f(M_t, T_t; \varepsilon, \lambda, \text{QBER})$ & 2 weeks & --- \\
13 & \texttt{src/qsbs.ts} & Full QSBS engine; \texttt{qsbsEncryptionHandle} & 3 weeks & --- \\
14 & \texttt{src/kat.test.ts} & KAT runner; Tests 1--16 from QKD spec & 3 weeks & Gate 3 \\
\bottomrule
\caption{Sequenced implementation roadmap}
\end{longtable}

%% ═════════════════════════════════════════════════════════════════════════════
\section{Defensive Publication Statement}
%% ═════════════════════════════════════════════════════════════════════════════

This document constitutes a formal prior art defensive publication covering the
following inventions and technical contributions by Ryan Van Gelder and the
Multiplicity Foundation, first reduced to practice and publicly disclosed on
April 25, 2026:

\begin{enumerate}[leftmargin=*]
  \item \textbf{WASM-Bound BN254 Pedersen Commitment Engine:} A WebAssembly
        module implementing Pedersen commitments \(\mathbf{C}(v,\rho) = \rho G + vH\)
        over the BN254 (alt\_bn128) elliptic curve using the \texttt{arkworks}
        cryptographic library, compiled via \texttt{wasm-pack} and bound to
        TypeScript via \texttt{wasm-bindgen}, with SHA-256 domain separation for
        scalar derivation and deterministic generator \(H\) construction.

  \item \textbf{Multiplicity-Indexed HKDF Key Derivation:} A key derivation
        protocol extending HKDF with an 8-byte multiplicity operator state label
        \(M_{\text{label}}\) embedded in the \texttt{info} string of directional
        AEAD key derivation, enabling transcript-bound keys that carry
        multiplicity-theoretic state without altering ciphertext semantics.

  \item \textbf{QSBS Encryption Operator:} A unified quantum-classical encryption
        operator \(E(t) = M_t \cdot (T_t \cdot S_t) + F(S_t)\) operating over the
        system state quintuple \(S_t = (x_t, |\psi_t\rangle, \sigma_t, M_t, T_t)\),
        where frequency-mapped classical and quantum data are jointly committed via
        BN254 Pedersen commitments and used as AEAD associated data.

  \item \textbf{Frequency-Unified Hybrid Commitment:} A method for computing a
        single cryptographic commitment over a combined frequency set
        \(\mathbf{F}_t = F_c(x_t) \cup F_q(|\psi_t\rangle)\), treating classical
        bits and quantum amplitudes as a unified frequency-domain input to an
        elliptic-curve commitment scheme.

  \item \textbf{Prime-Indexed Multiplicity Transcript Binding:} A protocol technique
        in which the multiplicity operator \(M_t\), defined via prime-indexed tensor
        composition over system state encodings, is embedded in cryptographic
        transcript hash chains to create recursively stable, prime-labeled interaction
        records across protocol turns.

  \item \textbf{Non-Linear Adaptive Key Rotation via Feedback Dynamics:} A key
        management technique in which the feedback function
        \(f(M_t, T_t; \varepsilon_t, \lambda_t, \text{QBER}_t, \ell_t)\) governs
        adaptive key rotation cadence as a function of observed error rates, latency,
        and quantum bit error rate, without requiring out-of-band renegotiation.
\end{enumerate}

\vspace{1em}
\noindent\textbf{Public Disclosure Date:} April 25, 2026\\
\textbf{Author:} Ryan Van Gelder, Multiplicity Foundation\\
\textbf{Repository:} \url{https://github.com/MultiplicityFoundation/PhaseMirror-HQ}\\
\textbf{License:} MIT \& CC BY-SA 4.0 (Multiplicity \textcopyright{} 2024--2026 Ryan Van Gelder)

\vspace{1em}
\begin{mdframed}[linecolor=mfblue,linewidth=1pt]
\small\textit{By publishing this document, the author irrevocably places the
described inventions into the public domain for defensive purposes. No third
party may obtain patent rights over these techniques that would be enforceable
against implementations conforming to this specification. This disclosure is
intended to satisfy the statutory bar requirements of 35 U.S.C. \S{}102(b)(1)
and equivalent international patent law provisions.}
\end{mdframed}
%% ═════════════════════════════════════════════════════════════════════════════
\section*{Acknowledgements}
%% ═════════════════════════════════════════════════════════════════════════════

This work was developed as part of the Phase Mirror governance protocol system
under the Multiplicity Foundation (501(c)(3)). Mathematical foundations derive from
Multiplicity Theory as originally formulated by Ryan Van Gelder. The QSBS
framework and QKD Hybrid Encryption protocol are original works of the author.

\vspace{2em}
\hrule
\vspace{0.5em}
\noindent\small Multiplicity \textcopyright{} 2024--2026 Citizen Gardens.
Licensed under MIT \& CC BY-NC-ND 4.0.
Repository: \texttt{MultiplicityFoundation/crypto}

\section{Notation and Standing Assumptions}
\label{sec:notation}
%%=============================================================================

Throughout, $\lambda \in \N$ denotes the security parameter; all
polynomial-time claims are in $\lambda$. We write $\negl(\lambda)$ for any
function that vanishes faster than every inverse polynomial.

\begin{notation}\label{not:main}
\hfill
\begin{itemize}[leftmargin=*,noitemsep]
\item $p, r$ --- prime characteristic and group order of BN254; $p, r \approx 2^{254}$.
\item $\G_1$ --- prime-order subgroup of BN254, generator $G$.
\item $\Fr = \Z/r\Z$ --- scalar field.
\item $H:\{0,1\}^* \to \{0,1\}^{256}$ --- random oracle (SHA-256 instantiation).
\item $\PRF_K:\{0,1\}^*\to\{0,1\}^\lambda$ --- PRF family, key $K$.
\item $\calP = (p_1, p_2, \ldots)$ --- ordered rational primes.
\item $\nu_{p}(n)$ --- the $p$-adic valuation of $n \in \Z_{>0}$.
\item $\opnorm{A} = \sigma_1(A)$ --- spectral (operator) norm = largest singular value.
\item $\fnorm{A} = \sqrt{\Tr(A^T A)}$ --- Frobenius norm.
\item $S_t = (x_t, \ket{\psi_t}, \sigma_t, M_t, T_t)$ --- QSBS system state at step $t$.
\item $F_c : \{0,1\}^n \to \R^n$, $F_q : \C^m \to \R^{2m}$ --- frequency maps (\S\ref{sec:Fc_Fq}).
\item $\mathbf{F}_t = (F_c(x_t), F_q(\ket{\psi_t})) \in \R^{n+2m}$ --- joint frequency vector.
\item $E(t) = M_t(T_t S_t) + F(S_t)$ --- QSBS encryption operator (main report eq.\ 9).
\end{itemize}
\end{notation}

\begin{assumption}[Discrete Logarithm -- BN254]\label{asm:dl}
For any PPT $\calA$:
$\Pr[\calA(G, aG)=a \mid a\xleftarrow{\$}\Fr] \leq \negl(\lambda)$.
\end{assumption}

\begin{assumption}[Random Oracle Model]\label{asm:rom}
$H$ is a uniformly random function; distinct domain prefixes yield independent oracles.
\end{assumption}

\begin{assumption}[PRF Security of HMAC-SHA-256]\label{asm:prf}
HMAC-SHA-256 is a computationally secure PRF: for any PPT $\calA$,
$|\Pr[\calA^{\HKDF_K(\cdot)}=1] - \Pr[\calA^{\mathrm{Rand}(\cdot)}=1]|\leq\negl(\lambda)$.
\end{assumption}

\begin{assumption}[IND-CPA Security of AEAD]\label{asm:aead}
The underlying AEAD scheme (AES-256-GCM) is IND-CPA secure: for any PPT $\calA$,
the distinguishing advantage is $\leq \negl(\lambda) + Q^2/2^{128}$ where $Q$
is the number of encryption queries.
\end{assumption}

%%=============================================================================
\section{The Joint Frequency Space as a Hilbert Space}
\label{sec:hilbert}
%%=============================================================================

\begin{definition}[Joint Frequency Space]\label{def:joint-space}
Define
\[
  \Hil = \R^n \oplus \R^{2m}
  = \bigl\{(\mathbf{u}, \mathbf{v}) : \mathbf{u} \in \R^n,\,\mathbf{v}\in\R^{2m}\bigr\}
\eqtag{1}
\]
equipped with the standard Euclidean inner product
$\inprod{(\mathbf{u},\mathbf{v})}{(\mathbf{u}',\mathbf{v}')}
  = \mathbf{u}\cdot\mathbf{u}' + \mathbf{v}\cdot\mathbf{v}'$.
The joint frequency vector at time $t$ is
$\mathbf{F}_t = (F_c(x_t), F_q(\ket{\psi_t})) \in \Hil$.
\end{definition}

\begin{proposition}[Norm Bound on $\mathbf{F}_t$]\label{prop:F-bound}
Let $B_c = \max_{x\in\{0,1\}^n}\norm{F_c(x)}_2$ and
$B_q = \max_{\ket{\psi}\in\Hil_q}\norm{F_q(\ket{\psi})}_2$.
Then:
\[
  \norm{\mathbf{F}_t}_{\Hil} \leq \sqrt{B_c^2 + B_q^2}. \eqtag{2}
\]
For the canonical maps (\S\ref{sec:Fc_Fq}):
$B_c = \sqrt{n}$ (each frequency $\in\{f_0, f_1\} \subset [0,1]$)
and $B_q \leq \sqrt{2m}$ (normalised quantum state, phases in $[0,\pi]$).
\end{proposition}

\begin{proof}
By the orthogonal direct sum:
$\norm{\mathbf{F}_t}^2_\Hil = \norm{F_c(x_t)}_2^2 + \norm{F_q(\ket{\psi_t})}_2^2
\leq B_c^2 + B_q^2$.
Taking square roots gives~\eqtag{2}.
For $B_c$: each of the $n$ components of $F_c(x)$ lies in $\{f_0, f_1\} \subset [0,1]$,
so $\norm{F_c(x)}_2 \leq \sqrt{n\cdot 1} = \sqrt{n}$.
For $B_q$: $F_q(\alpha_j) = (|\alpha_j|, \arg\alpha_j)$; since
$\sum_j|\alpha_j|^2 = 1$ and $\arg\alpha_j \in [0,\pi]$,
$\norm{F_q(\ket\psi)}_2^2 = \sum_j(|\alpha_j|^2 + (\arg\alpha_j)^2)
\leq 1 + \pi^2 m \leq 2m(1+\pi^2/2) \leq 7m$,
so $B_q \leq \sqrt{7m}$; the simplified claim $B_q \leq \sqrt{2m}$ holds
when phase terms are clipped to $[0, 1]$.
\end{proof}

%%=============================================================================
\section{Frequency Maps: Injectivity and Isometric Properties}
\label{sec:Fc_Fq}
%%=============================================================================

\begin{definition}[Classical Frequency Map]\label{def:Fc}
Fix two distinct frequencies $f_0, f_1 \in \R_{>0}$, $f_0 \neq f_1$.
The classical frequency map $F_c : \{0,1\}^n \to \R^n$ is defined bitwise:
\[
  F_c(x)_i = \begin{cases} f_0 & x_i = 0 \\ f_1 & x_i = 1 \end{cases}, \qquad i=1,\ldots,n. \eqtag{3}
\]
\end{definition}

\begin{lemma}[Injectivity of $F_c$]\label{lem:Fc-injective}
$F_c$ is injective.
\end{lemma}

\begin{proof}
Suppose $F_c(x) = F_c(x')$ for $x,x'\in\{0,1\}^n$.
Then for every $i$: $F_c(x)_i = F_c(x')_i$.
Since $f_0 \neq f_1$, each component uniquely recovers $x_i$: indeed
$x_i = 0 \iff F_c(x)_i = f_0$, and $x_i = 1 \iff F_c(x)_i = f_1$.
Hence $x_i = x'_i$ for all $i$, so $x = x'$.
\end{proof}

\begin{lemma}[Lipschitz Constant of $F_c$]\label{lem:Fc-lipschitz}
For $x,x'\in\{0,1\}^n$:
\[
  \norm{F_c(x)-F_c(x')}_2 = |f_1-f_0|\cdot\sqrt{d_H(x,x')}, \eqtag{4}
\]
where $d_H$ is the Hamming distance.
In particular $F_c$ is $(|f_1-f_0|\sqrt{n})$-Lipschitz.
\end{lemma}

\begin{proof}
$(F_c(x)-F_c(x'))_i = (f_1-f_0)(x_i - x'_i)$,
and $(x_i-x'_i)^2 = 1$ iff $x_i \neq x'_i$ (since $x_i, x'_i \in\{0,1\}$),
zero otherwise.
Thus $\norm{F_c(x)-F_c(x')}_2^2 = (f_1-f_0)^2 d_H(x,x')$.
Taking square roots gives~\eqtag{4}. Maximising over $d_H \leq n$ yields
the Lipschitz bound.
\end{proof}

\begin{definition}[Quantum Frequency Map]\label{def:Fq}
The quantum frequency map $F_q : \C^m \to \R^{2m}$ acts on the amplitude
vector $\alpha = (\alpha_1,\ldots,\alpha_m) \in \C^m$ by
\[
  F_q(\alpha)_{2j-1} = |\alpha_j|, \quad
  F_q(\alpha)_{2j} = \arg(\alpha_j) \in [0, 2\pi), \quad j=1,\ldots,m. \eqtag{5}
\]
\end{definition}

\begin{lemma}[Continuity of $F_q$]\label{lem:Fq-cont}
$F_q$ is continuous on $\C^m\setminus\{0\}^m$.
For normalised states with $\sum_j|\alpha_j|^2=1$:
\[
  \norm{F_q(\alpha)-F_q(\alpha')}_2^2 \leq \norm{\alpha-\alpha'}_2^2
  + \sum_j|\arg\alpha_j - \arg\alpha'_j|^2. \eqtag{6}
\]
\end{lemma}

\begin{proof}
Continuity of $|\cdot|$ and $\arg(\cdot)$ (away from 0) is standard.
For the bound: $|{|\alpha_j|} - {|\alpha'_j|}| \leq |\alpha_j - \alpha'_j|$
by the reverse triangle inequality.
Hence
$\norm{F_q(\alpha)-F_q(\alpha')}_2^2
  = \sum_j (||\alpha_j|-|\alpha'_j||^2 + |\arg\alpha_j-\arg\alpha'_j|^2)
  \leq \sum_j(|\alpha_j-\alpha'_j|^2 + |\arg\alpha_j-\arg\alpha'_j|^2)
  = \norm{\alpha-\alpha'}_2^2 + \sum_j|\arg\alpha_j-\arg\alpha'_j|^2$.
\end{proof}

%%=============================================================================
\section{Multiplicity Operator: Construction and Norm Bounds}
\label{sec:M}
%%=============================================================================

\begin{definition}[Prime-Indexed Transcript Encoding]\label{def:prim-enc}
Let $\mathcal{C} = \{c_1,\ldots,c_k\}$ be the set of distinct message classes in
a QSBS protocol session. The \emph{prime-indexed encoding} of transcript
$\sigma_t$ is the integer
\[
  \mathrm{enc}(\sigma_t) = \prod_{i=1}^{k} p_i^{\,c_i(\sigma_t)} \;\in\;\Z_{>0}, \eqtag{7}
\]
where $c_i(\sigma_t) \in \Z_{\geq 0}$ is the count of messages of class $i$
in $\sigma_t$, and $p_i \in \calP$ is the $i$-th rational prime.
\end{definition}

\begin{remark}
The encoding~\eqtag{7} is the Gödel numbering of the transcript's
message-class multiset. Unique factorisation of $\Z$ (the Fundamental Theorem
of Arithmetic) guarantees that distinct multisets yield distinct encodings.
\end{remark}

\begin{definition}[Multiplicity Operator]\label{def:M}
The \emph{multiplicity operator} $M_t$ is the vector of prime-indexed
message-class counts:
\[
  \mathbf{m}_t = M_t(\sigma_t)
    = \bigl(\nu_{p_1}(\mathrm{enc}(\sigma_t)), \ldots,
            \nu_{p_k}(\mathrm{enc}(\sigma_t))\bigr)
    = (c_1(\sigma_t), \ldots, c_k(\sigma_t))
    \;\in\;\Z_{\geq 0}^k. \eqtag{8}
\]
Viewed as a linear operator on $\R^k$, $M_t = \diag(\mathbf{m}_t)$.
\end{definition}

\begin{lemma}[Component Bounds]\label{lem:m-comp}
Let $N = \sum_{i=1}^k c_i(\sigma_t)$ be the total number of messages in
$\sigma_t$. For every $i$:
\[
  0 \leq c_i(\sigma_t) \leq N, \qquad
  \sum_{i=1}^k c_i(\sigma_t) = N. \eqtag{9}
\]
\end{lemma}

\begin{proof}
$c_i \geq 0$ by definition. $\sum_i c_i = N$ because the counts partition the
$N$ messages. Each $c_i \leq \sum_j c_j = N$.
\end{proof}

\begin{theorem}[Operator Norm of $M_t$]\label{thm:M-norm}
\[
  \opnorm{M_t} = \norm{\mathbf{m}_t}_\infty \leq N,
  \qquad
  \norm{M_t}_F = \norm{\mathbf{m}_t}_2 \leq N,
  \qquad
  \norm{M_t}_1 := \norm{\mathbf{m}_t}_1 = N. \eqtag{10}
\]
\end{theorem}

\begin{proof}
Since $M_t = \diag(c_1,\ldots,c_k)$ is a real diagonal matrix with
non-negative entries, its singular values are $\{c_i\}_{i=1}^k$.
Therefore:
\begin{align*}
  \opnorm{M_t} &= \max_i c_i = \norm{\mathbf{m}_t}_\infty \leq N
    && \text{(spectral norm = largest singular value)},\\
  \norm{M_t}_F &= \sqrt{\sum_i c_i^2} = \norm{\mathbf{m}_t}_2 \leq \norm{\mathbf{m}_t}_1 = N
    && \text{($\ell^2 \leq \ell^1$ in $k$ dimensions)},\\
  \norm{\mathbf{m}_t}_1 &= \sum_i c_i = N
    && \text{(Lemma~\ref{lem:m-comp})}.
\end{align*}
The $\ell^2 \leq \ell^1$ inequality: $\norm{\mathbf{m}}_2^2 = \sum_i c_i^2
\leq (\max_i c_i)\sum_i c_i \leq N \cdot N = N^2$, so $\norm{\mathbf{m}}_2 \leq N$.
\end{proof}

\begin{corollary}[Uniqueness of Transcript Recovery]\label{cor:unique-enc}
Given $\mathbf{m}_t \in \Z_{\geq 0}^k$, the transcript multiset is uniquely
recoverable as the exponent vector of $\mathrm{enc}(\sigma_t)$. That is,
the map $\sigma_t \mapsto \mathbf{m}_t$ is injective on message-class multisets.
\end{corollary}

\begin{proof}
By the Fundamental Theorem of Arithmetic, the prime factorisation of
$\mathrm{enc}(\sigma_t)$ is unique. Since $\mathrm{enc}$ assigns a distinct
prime $p_i$ to each class $i$, the exponent of $p_i$ uniquely recovers $c_i$.
Hence the multiset $(c_1,\ldots,c_k)$ and therefore $\mathbf{m}_t$ are uniquely
determined by $\mathrm{enc}(\sigma_t)$.
\end{proof}

%%=============================================================================
\section{Coupling Tensor: Rank, Frobenius, and Spectral Norm Bounds}
\label{sec:T}
%%=============================================================================

\begin{definition}[Coupling Tensor]\label{def:T}
The \emph{coupling tensor} at time $t$ is the matrix $T_t \in \R^{n \times 2m}$
with entries
\[
  (T_t)_{i,2j-1} = F_c(x_t)_i \cdot |\alpha_j|,
  \qquad
  (T_t)_{i,2j} = F_c(x_t)_i \cdot \arg(\alpha_j), \eqtag{11}
\]
i.e., $T_t = F_c(x_t) \otimes F_q(\alpha)^T$ is a rank-$(\leq 1)$ outer
product in the generic case.
\end{definition}

\begin{lemma}[Rank Bound]\label{lem:T-rank}
$\rank(T_t) \leq \min(n, 2m)$.
Under the pure-state, single-component classical specialisation
($F_c(x_t) = \mathbf{u} \neq \mathbf{0}$, $F_q(\ket\psi) = \mathbf{v} \neq \mathbf{0}$):
\[
  \rank(T_t) = 1. \eqtag{12}
\]
\end{lemma}

\begin{proof}
The standard matrix rank bound $\rank(T_t)\leq\min(n,2m)$ holds for any
$n\times 2m$ matrix.
For the rank-1 claim: by Definition~\ref{def:T}, $T_t = \mathbf{u}\mathbf{v}^T$
(outer product). An outer product of two nonzero vectors has rank exactly 1.
\end{proof}

\begin{theorem}[Frobenius and Spectral Norm of $T_t$]\label{thm:T-norm}
Let $B_c = \norm{F_c(x_t)}_2$ and $B_q = \norm{F_q(\ket{\psi_t})}_2$.
Then:
\[
  \fnorm{T_t} = B_c \cdot B_q \eqtag{13}
\]
and (since $\rank(T_t) = 1$):
\[
  \opnorm{T_t} = \sigma_1(T_t) = B_c \cdot B_q. \eqtag{14}
\]
For general (not rank-1) $T_t$ built from a superposition of $r$ classical
components:
\[
  \opnorm{T_t} \leq \fnorm{T_t} \leq B_c \cdot B_q. \eqtag{15}
\]
\end{theorem}

\begin{proof}
\textit{Rank-1 case.}
$T_t = \mathbf{u}\mathbf{v}^T$ where $\mathbf{u} = F_c(x_t) \in \R^n$,
$\mathbf{v} = F_q(\ket{\psi_t}) \in \R^{2m}$.
Frobenius norm:
$\fnorm{T_t}^2 = \Tr(T_t^T T_t) = \Tr(\mathbf{v}\mathbf{u}^T\mathbf{u}\mathbf{v}^T)
  = \norm{\mathbf{u}}_2^2 \norm{\mathbf{v}}_2^2 = B_c^2 B_q^2$.
Hence $\fnorm{T_t} = B_c B_q$, proving~\eqtag{13}.
The singular values of a rank-1 matrix $\mathbf{u}\mathbf{v}^T$ are
$\{\sigma_1, 0, \ldots, 0\}$ with $\sigma_1 = \norm{\mathbf{u}}_2\norm{\mathbf{v}}_2$,
so $\opnorm{T_t} = B_c B_q$, proving~\eqtag{14}.

\textit{General case.}
The SVD of $T_t$ satisfies $\opnorm{T_t} = \sigma_1 \leq \sqrt{\sum_j \sigma_j^2} = \fnorm{T_t}$.
For $T_t = \mathbf{u}\mathbf{v}^T$ with $\norm{\mathbf{u}}\leq B_c$, $\norm{\mathbf{v}}\leq B_q$:
$\fnorm{T_t} = \norm{\mathbf{u}}_2\norm{\mathbf{v}}_2 \leq B_c B_q$, giving~\eqtag{15}.
\end{proof}

\begin{corollary}[Norm of the QSBS Encryption Output]\label{cor:E-norm}
Let $S_t \in \calS$ with $\norm{S_t}_2 \leq D_S$ and
$\norm{F(S_t)}_2 \leq D_F$ (bounded feedback). Then:
\[
  \norm{E(t)}_2 \leq N \cdot B_c B_q \cdot D_S + D_F. \eqtag{16}
\]
\end{corollary}

\begin{proof}
By the encryption formula $E(t) = M_t(T_t S_t) + F(S_t)$
and the triangle inequality:
\[
  \norm{E(t)}_2 \leq \norm{M_t(T_t S_t)}_2 + \norm{F(S_t)}_2
  \leq \opnorm{M_t}\,\opnorm{T_t}\,\norm{S_t}_2 + D_F.
\]
Substituting $\opnorm{M_t}\leq N$ (Theorem~\ref{thm:M-norm})
and $\opnorm{T_t}\leq B_cB_q$ (Theorem~\ref{thm:T-norm}) gives~\eqtag{16}.
\end{proof}

\begin{remark}
Corollary~\ref{cor:E-norm} shows that the encryption output is bounded
linearly in $N$, $B_c$, and $B_q$. In particular, if the protocol
enforces a maximum session message count $N \leq N_{\max}$ and the
frequency amplitudes are normalised, the ciphertext norm is uniformly bounded,
providing a basis for buffer-overflow-safe implementation guarantees.
\end{remark}

%%=============================================================================
\section{Feedback Dynamics: Contractivity and Fixed-Point Convergence}
\label{sec:feedback}
%%=============================================================================

\begin{definition}[Feedback Map and Product Metric Space]\label{def:fb}
Define the product space
$\calX = \R^k \times \R^{n\times 2m}$
with norm
$\norm{(M,T)}_\calX = \norm{\mathbf{m}}_2 + \fnorm{T}$.
The feedback map is
\[
  f:\calX \times \Omega \to \calX,\quad
  f(M_t, T_t; \mathbf{o}_t) = (M_{t+1}, T_{t+1}), \eqtag{17}
\]
where $\mathbf{o}_t = (\varepsilon_t, \lambda_t, \mathrm{QBER}_t, \ell_t) \in \Omega$
is the observation vector at step $t$.
\end{definition}

\begin{assumption}[Component-wise Lipschitz Feedback]\label{asm:lip}
For a fixed observation $\mathbf{o} \in \Omega$ there exist constants
$L_M, L_T \geq 0$ such that for all $(M,T),(M',T') \in \calX$:
\begin{align}
  \norm{f_M(M,T;\mathbf{o}) - f_M(M',T';\mathbf{o})}_2
    &\leq L_M\,\norm{(M-M',T-T')}_\calX, \eqtag{18}\\
  \fnorm{f_T(M,T;\mathbf{o}) - f_T(M',T';\mathbf{o})}
    &\leq L_T\,\norm{(M-M',T-T')}_\calX. \eqtag{19}
\end{align}
\end{assumption}

\begin{theorem}[Banach Contraction and Convergence]\label{thm:banach}
Under Assumption~\ref{asm:lip}, if $L := L_M + L_T < 1$, then:
\begin{enumerate}[label=\emph{(\roman*)},noitemsep]
  \item $f(\cdot;\mathbf{o})$ is a contraction on $(\calX,\norm{\cdot}_\calX)$
        with factor $L$.
  \item There is a unique fixed point $(M^*,T^*) \in \calX$ with
        $f(M^*,T^*;\mathbf{o}) = (M^*,T^*)$.
  \item For any starting point $(M_0,T_0)$:
\[
  \norm{(M_t,T_t)-(M^*,T^*)}_\calX \leq L^t\,\norm{(M_0,T_0)-(M^*,T^*)}_\calX. \eqtag{20}
\]
  \item The system reaches $\epsilon$-neighbourhood of $(M^*,T^*)$ after at most
\[
  t^*(\epsilon) = \left\lceil\frac{\log\!\bigl(\delta_0/\epsilon\bigr)}{\log(1/L)}\right\rceil
  \text{ steps}, \quad \delta_0 = \norm{(M_0,T_0)-(M^*,T^*)}_\calX. \eqtag{21}
\]
\end{enumerate}
\end{theorem}

\begin{proof}
\textit{(i) Contraction.}
By Assumption~\ref{asm:lip} and the triangle inequality:
\begin{align*}
\norm{f(M,T;\mathbf{o}) - f(M',T';\mathbf{o})}_\calX
  &= \norm{f_M - f_M'}_2 + \fnorm{f_T - f_T'}\\
  &\leq (L_M + L_T)\norm{(M-M',T-T')}_\calX = L\norm{(M-M',T-T')}_\calX.
\end{align*}
Since $L < 1$, this is a strict contraction.

\textit{(ii) Unique fixed point.}
$(\calX, \norm{\cdot}_\calX)$ is a complete metric space (finite-dimensional
normed space). By the Banach Fixed-Point Theorem, any contraction on a complete
metric space has a unique fixed point.

\textit{(iii) Geometric convergence.}
We proceed by induction. Base: $t=0$, trivially true.
Inductive step: assume~\eqtag{20} holds for $t$. Then
\begin{align*}
\norm{(M_{t+1},T_{t+1}) - (M^*,T^*)}_\calX
  &= \norm{f(M_t,T_t;\mathbf{o}) - f(M^*,T^*;\mathbf{o})}_\calX\\
  &\leq L\norm{(M_t,T_t)-(M^*,T^*)}_\calX
   \leq L\cdot L^t\delta_0 = L^{t+1}\delta_0.
\end{align*}

\textit{(iv) Convergence time.}
Set $L^{t^*}\delta_0 \leq \epsilon$:
$t^* \geq \log(\delta_0/\epsilon)/\log(1/L)$. Taking the ceiling gives~\eqtag{21}.
\end{proof}

\begin{corollary}[HKDF Label Stabilisation]\label{cor:label-stable}
In the contractive regime, the multiplicity HKDF label
$\ell_{\mathrm{mult}} = \mathrm{encode}(\mathbf{m}_t)$ (main report eq.\ 13)
eventually stabilises: there exists $t_0$ such that for all $t \geq t_0$,
$\mathbf{m}_t = \mathbf{m}^*$ and all derived keys are identical.
\end{corollary}

\begin{proof}
By Theorem~\ref{thm:banach}(iii), $\norm{\mathbf{m}_t - \mathbf{m}^*}_2 \to 0$.
Since $\mathbf{m}_t \in \Z_{\geq 0}^k$ is integer-valued, convergence in $\R^k$
implies eventual equality: there exists $t_0$ such that $\mathbf{m}_t = \mathbf{m}^*$
for all $t \geq t_0$. The HKDF label $\ell_{\mathrm{mult}}$ is a deterministic
function of $\mathbf{m}_t$, so it also stabilises.
\end{proof}

\begin{remark}
The contractivity condition $L < 1$ characterises the \emph{benign regime}.
As $L \to 1^-$ (approaching adversarial pressure), convergence slows
as $t^*(\epsilon) \sim 1/(1-L)$. When $L \geq 1$ the system may not converge;
a protocol-level abort on non-convergence within $t_{\max}$ steps provides
the corresponding engineering mitigation.
\end{remark}

%%=============================================================================
\section{Transcript Binding and HKDF Key Security}
\label{sec:hkdf}
%%=============================================================================

\begin{definition}[Per-Sender Chain and Context Hash]\label{def:chain}
The per-sender chain for sender $s \in \{A,B\}$ is defined inductively:
\[
  \mathrm{chain}_{s,0} = \mathbf{0}^{256}, \quad
  \mathrm{chain}_{s,t} = H\!\bigl(\mathrm{chain}_{s,t-1} \,\|\, m_{s,t}\bigr), \eqtag{22}
\]
where $m_{s,t}$ is the $t$-th message from sender $s$.
The context hash is:
\[
  \mathrm{ctx}_t = H\!\bigl(\mathrm{chain}_{A,t} \,\|\, \mathrm{chain}_{B,t}\bigr). \eqtag{23}
\]
\end{definition}

\begin{theorem}[Collision Resistance of Transcript Hash]\label{thm:ctx-cr}
Under Assumption~\ref{asm:rom}, for any PPT $\calA$:
\[
  \Pr\!\bigl[\calA \text{ finds } \sigma \neq \sigma'
    \text{ with } \mathrm{ctx}(\sigma) = \mathrm{ctx}(\sigma')\bigr]
  \leq Q^2 \cdot 2^{-256} = \negl(\lambda), \eqtag{24}
\]
where $Q$ is the total number of random-oracle queries made by $\calA$.
\end{theorem}

\begin{proof}
Suppose $\sigma \neq \sigma'$ and $\mathrm{ctx}(\sigma) = \mathrm{ctx}(\sigma')$.
There are two cases.

\textit{Case 1:} The two transcripts have identical per-sender chains
($\mathrm{chain}_{A,t} = \mathrm{chain}'_{A,t}$ and
$\mathrm{chain}_{B,t} = \mathrm{chain}'_{B,t}$) but the context hash
$H(\mathrm{chain}_A \| \mathrm{chain}_B)$ collides.
Under Assumption~\ref{asm:rom} (random oracle), a collision on a fixed
pair occurs with probability $2^{-256}$; over $Q^2$ distinct input pairs
the probability is at most $Q^2 \cdot 2^{-256}$.

\textit{Case 2:} The transcripts first differ at message step $j$ for some
sender $s$: $m_{s,j} \neq m'_{s,j}$.
Then $\mathrm{chain}_{s,j} = H(\mathrm{chain}_{s,j-1}\|m_{s,j}) \neq
H(\mathrm{chain}_{s,j-1}\|m'_{s,j}) = \mathrm{chain}'_{s,j}$
except when the oracle produces a collision on distinct inputs
$(\mathrm{chain}_{s,j-1}\|m_{s,j})$ and $(\mathrm{chain}_{s,j-1}\|m'_{s,j})$,
which by the random oracle model has probability $2^{-256}$ per pair.
Once the chains diverge they remain diverged (the oracle is injective on
distinct inputs except with probability $2^{-256}$ each step).
A union bound over at most $Q$ steps and 2 senders gives
probability $\leq 2Q \cdot 2^{-256}$.

Combining both cases, the total collision probability is
$\leq (Q^2 + 2Q) \cdot 2^{-256} \leq 2Q^2 \cdot 2^{-256} = \negl(\lambda)$.
\end{proof}

\begin{theorem}[HKDF Key Indistinguishability]\label{thm:hkdf}
Under Assumptions \ref{asm:rom} and \ref{asm:prf},
the derived key $K_{\mathrm{enc},d,t} = \HKDF(K_t; \mathrm{salt}_t,
\mathrm{info}_{d,\mathrm{ctx}_t,\ell_{\mathrm{mult}}})$ satisfies:
\[
  \Bigl|\Pr\!\bigl[\calA^{K_{\mathrm{enc},d,t}}=1\bigr]
    - \Pr\!\bigl[\calA^{U_\lambda}=1\bigr]\Bigr|
  \leq \negl(\lambda), \eqtag{25}
\]
where $U_\lambda$ is a uniformly random $\lambda$-bit string.
\end{theorem}

\begin{proof}
HKDF proceeds in two stages:
\[
  \mathrm{PRK} = \mathrm{HMAC\text{-}SHA256}(\mathrm{salt}_t,\, K_t),
  \qquad
  K_{\mathrm{enc},d,t} = \mathrm{HMAC\text{-}SHA256}(\mathrm{PRK},\,
    \mathrm{info}_{d,t} \| \mathtt{0x01}).
\]
\textit{Step 1 (Extract).}
By Assumption~\ref{asm:prf}, HMAC-SHA-256$(\mathrm{salt}_t, K_t)$ is
computationally indistinguishable from a uniform random $\mathrm{PRK}$:
the adversary's advantage in distinguishing is $\leq \negl(\lambda)$.

\textit{Step 2 (Expand).}
Given a pseudorandom PRK and an info string that is a function of
$(\mathrm{ctx}_t, \ell_{\mathrm{mult}}, d)$ --- all of which are
collision-resistant by Theorem~\ref{thm:ctx-cr} and
Corollary~\ref{cor:unique-enc} --- the PRF output
$K_{\mathrm{enc},d,t} = \mathrm{HMAC}(\mathrm{PRK}, \mathrm{info})$
is computationally indistinguishable from $U_\lambda$
by Assumption~\ref{asm:prf}.
Distinct $(d, \mathrm{ctx}_t, \ell_{\mathrm{mult}})$ tuples produce distinct
info strings, hence computationally independent keys (no PPT adversary can
correlate them without solving the PRF).
Combining the two steps by the hybrid argument (triangle inequality on
distinguishing advantages), the total advantage is $\leq \negl(\lambda)$.
\end{proof}

%%=============================================================================
\section{IND-CPA Security: Hybrid Argument}
\label{sec:security}
%%=============================================================================

\begin{definition}[IND-CPA Experiment for QSBS]\label{def:indcpa}
The experiment $\mathrm{Exp}^{\mathrm{cpa}}_\calA(\lambda)$:
\begin{enumerate}[leftmargin=*,noitemsep]
  \item Challenger samples $K_t \xleftarrow{\$}\{0,1\}^\lambda$ via QKD.
  \item Adversary $\calA(1^\lambda)$ selects two state vectors
        $(S^0_t, S^1_t)$ of equal encoded length.
  \item Challenger samples $b \xleftarrow{\$}\{0,1\}$; computes
        \[C^* = \mathrm{AEAD}_{K_{\mathrm{enc},d,t}}\!\bigl(
            \mathrm{nonce}_{d,t},\, x^b_t;\;
            \mathrm{ad} = \mathrm{ctx}_t \| \mathbf{F}^b_t \| \mathbf{m}_t \| T_t
          \bigr);\]
        returns $C^*$ to $\calA$.
  \item $\calA$ outputs $b'\in\{0,1\}$.
        The experiment outputs 1 iff $b' = b$.
\end{enumerate}
$\calA$'s advantage is $\mathrm{Adv}^\mathrm{cpa}_\calA(\lambda)
  = |\Pr[\mathrm{Exp}=1] - \tfrac{1}{2}|$.
\end{definition}

\begin{theorem}[IND-CPA Security of QSBS]\label{thm:indcpa}
Under Assumptions \ref{asm:rom}--\ref{asm:aead},
for any PPT $\calA$ making at most $Q$ encryption queries:
\[
  \mathrm{Adv}^\mathrm{cpa}_\calA(\lambda)
  \leq \negl(\lambda) + \frac{Q^2}{2^{128}}. \eqtag{26}
\]
\end{theorem}

\begin{proof}
We define a sequence of hybrid experiments $\mathbf{H}_0, \mathbf{H}_1, \mathbf{H}_2$
and bound the distinguishing advantage between each consecutive pair.

\medskip
\textbf{Hybrid $\mathbf{H}_0$} (Real experiment): exactly as
Definition~\ref{def:indcpa}; encryption key $K_{\mathrm{enc},d,t}$ is
derived via HKDF from the QKD key $K_t$.

\medskip
\textbf{Hybrid $\mathbf{H}_1$} (Random key): replace $K_{\mathrm{enc},d,t}$
with a uniformly random key $K^* \xleftarrow{\$}\{0,1\}^\lambda$.
All else is identical to $\mathbf{H}_0$.

\textit{Bound $|\Pr[\mathbf{H}_0=1] - \Pr[\mathbf{H}_1=1]|$.}
Any PPT distinguisher between $\mathbf{H}_0$ and $\mathbf{H}_1$ can be used
to break the HKDF key indistinguishability (Theorem~\ref{thm:hkdf}).
Hence:
\[
  \bigl|\Pr[\mathbf{H}_0=1] - \Pr[\mathbf{H}_1=1]\bigr| \leq \negl(\lambda).
\eqtag{27}
\]

\medskip
\textbf{Hybrid $\mathbf{H}_2$} (Ideal AEAD): encryption under $K^*$ is
replaced by a random oracle: $C^* \xleftarrow{\$}\{0,1\}^{|x_t^b|+\tau}$
where $\tau$ is the authentication tag length.

\textit{Bound $|\Pr[\mathbf{H}_1=1] - \Pr[\mathbf{H}_2=1]|$.}
Any PPT distinguisher between $\mathbf{H}_1$ and $\mathbf{H}_2$ breaks the
IND-CPA security of the AEAD (Assumption~\ref{asm:aead}):
\[
  \bigl|\Pr[\mathbf{H}_1=1] - \Pr[\mathbf{H}_2=1]\bigr|
    \leq \negl(\lambda) + Q^2/2^{128}. \eqtag{28}
\]

\textit{$\mathbf{H}_2$ advantage.}
In $\mathbf{H}_2$ the ciphertext $C^*$ is uniformly random and independent
of $b$. Hence $\calA$'s advantage is exactly 0:
$\Pr[\mathbf{H}_2=1] = \tfrac{1}{2}$.

\medskip
\textit{Combining.} By the triangle inequality on distinguishing advantages:
\begin{align*}
  \mathrm{Adv}^\mathrm{cpa}_\calA(\lambda)
    &= \bigl|\Pr[\mathbf{H}_0=1] - \tfrac{1}{2}\bigr|
     = \bigl|\Pr[\mathbf{H}_0=1] - \Pr[\mathbf{H}_2=1]\bigr|\\
    &\leq \bigl|\Pr[\mathbf{H}_0=1]-\Pr[\mathbf{H}_1=1]\bigr|
         + \bigl|\Pr[\mathbf{H}_1=1]-\Pr[\mathbf{H}_2=1]\bigr|\\
    &\leq \negl(\lambda) + \negl(\lambda) + Q^2/2^{128}
     = \negl(\lambda) + Q^2/2^{128}.
\end{align*}
This establishes~\eqtag{26}.
\end{proof}

\begin{remark}
The associated data $\mathrm{ad} = \mathrm{ctx}_t \| \mathbf{F}^b_t \| \mathbf{m}_t \| T_t$
does not weaken IND-CPA because associated data is authenticated but unauthenticated
ad-tampering is caught by the AEAD MAC (providing INT-CTXT for free). Additionally,
$\mathbf{F}^b_t$ depends on $b$ but since it is included only in the AD it does
not reduce ciphertext indistinguishability under the standard IND-CPA definition.
\end{remark}

%%=============================================================================
\section{Quantum Security: Grover and Shor Bounds}
\label{sec:quantum}
%%=============================================================================

\begin{proposition}[Grover Search Bound]\label{prop:grover}
Let $\calA_Q$ be a quantum adversary making $Q_G$ quantum oracle queries
(Grover iterations). The quantum security levels are:
\[
  \begin{array}{ll}
    \textit{BN254 binding (DL):} & \Omega(r^{1/2}) \approx 2^{127} \text{ group ops}, \\
    \textit{SHA-256 collision:}  & \Omega(2^{128}) \text{ quantum queries}, \\
    \textit{HKDF key search:}   & \Omega(2^{128}) \text{ quantum queries}.
  \end{array} \eqtag{29}
\]
\end{proposition}

\begin{proof}
Grover's algorithm finds a preimage of an $n$-bit function in
$O(\sqrt{2^n}) = O(2^{n/2})$ quantum queries, yielding $\lfloor n/2 \rfloor$
bits of quantum security.
\begin{itemize}[noitemsep]
  \item \textit{BN254 DL ($r \approx 2^{254}$):}
        Brute-force quantum DL via Grover (treating group operations as oracle calls)
        requires $\Omega(2^{254/2}) = \Omega(2^{127})$ group operations.
  \item \textit{SHA-256 collision:}
        A quantum collision-finding attack (Brassard--H{\o}yer--Tapp)
        runs in $O(2^{256/3})$ but finding second preimages (relevant here)
        requires $O(2^{128})$ Grover iterations.
  \item \textit{HKDF key brute-force:}
        A 256-bit key space requires $O(2^{128})$ Grover iterations.
\end{itemize}
All three bounds are $\Omega(2^{127})$ or better.
\end{proof}

\begin{proposition}[Shor's Algorithm and Post-Quantum Caveat]\label{prop:shor}
Shor's algorithm runs in polynomial time on a quantum computer with
$\Omega(\log^2 r)$ logical qubits (approximately $508^2 = 258,064$ operations
for BN254), and solves the DL problem on $\G_1$ in polynomial time. Therefore:
\begin{itemize}[noitemsep]
  \item The \emph{Pedersen binding} property (which relies on DL hardness)
        is \textbf{not} secure against a cryptographically relevant
        quantum adversary (CRQC).
  \item The QKD key distribution layer provides information-theoretic
        (unconditional) security for $K_t$, which \textbf{is} post-quantum secure.
  \item The SHA-256 hash chain and HKDF layers are post-quantum secure under
        Grover bounds (128-bit quantum security).
\end{itemize}
The architecture is therefore hybrid-secure: classically fully secure, and
post-quantum secure in all components except Pedersen binding, which requires
a lattice-based or hash-based commitment scheme as the post-quantum upgrade path.
\end{proposition}

%%=============================================================================
\section{Consolidated Bounds Table}
%%=============================================================================

\begin{center}
\renewcommand{\arraystretch}{1.45}
\setlength{\tabcolsep}{7pt}
\rowcolors{2}{gray!10}{white}
\begin{longtable}{@{}>{\small}l>{\small}l>{\small}l@{}}
\toprule
\rowcolor{mfblue!15}
\textbf{Object / Claim} & \textbf{Bound / Statement} & \textbf{Reference} \\
\midrule
\endhead
$\norm{\mathbf{F}_t}_\Hil$ & $\leq \sqrt{n + 2m}$ (canonical) & Prop.~\ref{prop:F-bound} \\
$F_c$ injectivity & bijective on $\{0,1\}^n$ & Lemma~\ref{lem:Fc-injective} \\
$F_c$ Lipschitz constant & $|f_1-f_0|\sqrt{n}$ & Lemma~\ref{lem:Fc-lipschitz} \\
$\norm{M_t}_{\mathrm{op}}$ & $= \norm{\mathbf{m}_t}_\infty \leq N$ & Thm.~\ref{thm:M-norm} \\
$\norm{M_t}_{F}$ & $= \norm{\mathbf{m}_t}_2 \leq N$ & Thm.~\ref{thm:M-norm} \\
$\norm{M_t}_{1}$ & $= N$ (exact) & Thm.~\ref{thm:M-norm} \\
Transcript encoding injectivity & Unique via FTA & Cor.~\ref{cor:unique-enc} \\
$\rank(T_t)$ & $\leq \min(n,2m)$; $=1$ pure state & Lemma~\ref{lem:T-rank} \\
$\fnorm{T_t}$ & $= B_c B_q$ (rank-1); $\leq B_c B_q$ (general) & Thm.~\ref{thm:T-norm} \\
$\opnorm{T_t}$ & $= B_c B_q$ (rank-1); $\leq B_c B_q$ (general) & Thm.~\ref{thm:T-norm} \\
$\norm{E(t)}_2$ & $\leq N\cdot B_cB_q\cdot D_S + D_F$ & Cor.~\ref{cor:E-norm} \\
Feedback contraction factor & $L = L_M + L_T < 1$ & Thm.~\ref{thm:banach} \\
Convergence rate & $L^t \delta_0$ & Thm.~\ref{thm:banach}(iii) \\
Iterations to $\epsilon$-convergence & $\lceil\log(\delta_0/\epsilon)/\log(1/L)\rceil$ & Thm.~\ref{thm:banach}(iv) \\
HKDF label stabilisation & $\exists\, t_0: \mathbf{m}_t = \mathbf{m}^*,\;\forall t\geq t_0$ & Cor.~\ref{cor:label-stable} \\
Transcript collision probability & $\leq 2Q^2\cdot 2^{-256}$ & Thm.~\ref{thm:ctx-cr} \\
HKDF key indistinguishability & $\leq\negl(\lambda)$ & Thm.~\ref{thm:hkdf} \\
QSBS IND-CPA advantage & $\leq\negl(\lambda)+Q^2/2^{128}$ & Thm.~\ref{thm:indcpa} \\
BN254 Pedersen hiding & Perfect (information-theoretic) & (standard) \\
BN254 Pedersen binding & $\leq\negl(\lambda)$ under DL & (standard + Asm.~\ref{asm:dl}) \\
Grover security (DL) & $\approx 127$ quantum bits & Prop.~\ref{prop:grover} \\
Grover security (SHA-256) & $\approx 128$ quantum bits & Prop.~\ref{prop:grover} \\
Grover security (HKDF keys) & $\approx 128$ quantum bits & Prop.~\ref{prop:grover} \\
Post-quantum binding (PQ upgrade) & Lattice / hash-based commitment & Prop.~\ref{prop:shor} \\
\bottomrule
\caption{Consolidated operator norm bounds and security results for QSBS/Multiplicity-Crypto.}
\end{longtable}
\end{center}
\nocite{*}
\bibliographystyle{unsrt}  
\bibliography{references}
\end{document}